\documentclass[11pt]{amsart}
\usepackage[T1]{fontenc}
\usepackage{lmodern}
\usepackage{amsmath,amssymb}
\usepackage{microtype}
\usepackage[hidelinks,pdftitle={Draft: A matroid proof of bipartite contractibility}]{hyperref}
\setlength{\emergencystretch}{1em}
\raggedbottom

\newtheorem{theorem}{Theorem}[section]
\newtheorem{lemma}[theorem]{Lemma}
\newtheorem{corollary}[theorem]{Corollary}

\title[Bipartite contractibility]{A matroid proof of bipartite contractibility}
\date{Draft, 8 September 2026}

\begin{document}
\begin{abstract}
We give an independent proof that every scheme of a finite bipartite
graph contains a minor rooted at all prescribed vertices. This supplies
the intended conclusion of an earlier bipartite-flow assertion of
Biswal, Lee and Rao, whose prefix construction fails. The proof uses
matroid union to allocate disjoint forests: either spanning trees give
the model, or simultaneous component contractions yield a smaller
scheme of the same target while preserving every root.
\end{abstract}
\maketitle

\section{Statement and colour normalisation}

All graphs are finite. Targets and hosts are simple; auxiliary
multigraphs may have parallel edges. Following K\"undgen, Pelsmajer and
Ramamurthi~\cite[Definitions 1.1 and 2.1]{KPR}, an \emph{$H$-scheme}
in a graph $G$ consists of distinct roots $(r_v:v\in V(H))$ and a simple
$r_u,r_v$-path $P_{uv}$ for every $uv\in E(H)$, with no other root
internal to a path, such that any nonempty collection of paths with a common
vertex has a common target endpoint.
A \emph{rooted $H$-minor model} is a family of pairwise disjoint nonempty
connected vertex sets $(C_v:v\in V(H))$, with $r_v\in C_v$, such that
$C_u,C_v$ are adjacent whenever $uv\in E(H)$. A graph $H$ is
\emph{contractible} if every $H$-scheme contains such a model.

\begin{theorem}\label{thm:main}
Every finite simple bipartite graph is contractible.
\end{theorem}

This answers the bipartite contractibility questions
in~\cite[Section 8, Questions 2--4]{KPR}, including $K_{2,4}$, $K_{3,3}$
and bipartite theta graphs. The earlier flow assertion is discussed in
Section~\ref{sec:flow}.
The proof uses Edmonds' matroid union rank formula and the following
part of the coloured-scheme reduction of~\cite[Lemma 3.3]{KPR}.
A minor is \emph{root-preserving} when its vertices have
pairwise disjoint connected preimages and each new root's preimage
contains its corresponding original root.

\begin{lemma}\label{lem:normalization}
Every $H$-scheme has a root-preserving minor whose underlying graph is
properly coloured by a map $f$ to $V(H)$, with $f(r_v)=v$, and whose
scheme paths use only their endpoint colours.
\end{lemma}

\begin{proof}
Delete unused edges and nonroot vertices. Give each root its target
label. At each nonroot, choose an endpoint common to all target edges
whose paths contain it. Thus each path uses only its endpoint labels.
Contract every connected component induced by one label. Such a component
contains at most one root; the quotient colouring is proper. Each path
maps to an endpoint-coloured walk, from which deleting closed excursions
gives a simple path between the new roots. Another root cannot occur,
since it has a different label. All paths meeting at a vertex have its
label as a common target endpoint. Delete unused vertices and edges,
retaining isolated prescribed roots.
\end{proof}

Contractibility is equivalent to property $(*)$ of Kriesell and
Mohr~\cite[Definition 1]{KM}. There the roots form a transversal, one
vertex from each class of a proper colouring, and each required pair is
joined by a path using its two colours. Such paths form a scheme; conversely,
Lemma~\ref{lem:normalization} reduces every scheme to this setting and
the resulting rooted model lifts. Thus every bipartite graph has
property $(*)$.

\section{A matroid union equality}

Let $(M_j:j\in J)$ be a finite family of multigraphs, each with distinct
edge labels in a subset $E_j$ of a finite set $E$. Extend its graphic
matroid to $E$ by declaring labels outside $E_j$ to be matroid loops.
Write $r_j$ for its rank function and $M_j(X)$ for its spanning subgraph
with labels in $X$. If $c_j(X)$ counts its components, including isolated
vertices, then $r_j(X)=|V(M_j)|-c_j(X)$.

\begin{lemma}\label{lem:union}
Let $R$ be the maximum of $\sum_j|I_j|$ over pairwise disjoint sets
$I_j\subseteq E_j$ that are forests in their respective graphs.
There is $X\subseteq E$ such that
\begin{equation}\label{eq:union}
 R=|E\setminus X|+\sum_jr_j(X).
\end{equation}
For any maximising family $(I_j)$ and any such $X$, the forest
$I_j\cap X$ spans every component of $M_j(X)$.
\end{lemma}

\begin{proof}
Edmonds' matroid union rank formula~\cite[Theorem 1 and the following
discussion, pp.~202--203]{Edmonds} gives
\[
 R=\min_{X\subseteq E}\left(|E\setminus X|+\sum_jr_j(X)\right).
\]
For a minimising $X$ and a maximising disjoint family,
\[
 R=\sum_j|I_j\setminus X|+\sum_j|I_j\cap X|
 \leq |E\setminus X|+\sum_jr_j(X)=R.
\]
Equality forces $|I_j\cap X|=r_j(X)$ for every $j$. An independent
set of this rank is a spanning tree in each nontrivial component of
$M_j(X)$, with isolated vertices retained.
\end{proof}

\section{A simultaneous component contraction}

Let $H$ be bipartite with parts $(A,B)$, and let $G$ be the union of a
properly endpoint-coloured $H$-scheme and its isolated roots, as in
Lemma~\ref{lem:normalization}. Put
\[
 W_a=f^{-1}(a)\quad(a\in A),\qquad
 E=f^{-1}(B)\setminus\{r_b:b\in B\}.
\]
For $a\in A$, construct a multigraph $M_a$ on $W_a$. Each actual
occurrence of $x\in E$ on $P_{ab}$ gives an edge labelled $x$ between
its two neighbours on that path. They are distinct vertices of colour
$a$, so the edge is not a loop. Each label occurs at most once in
$M_a$, since its colour determines $b$ and the target is simple.
Denote the label set of $M_a$ by $E_a$.

Suppressing the internal $B$ vertices of $P_{ab}$ and deleting its
terminal root $r_b$ gives a path in $M_a$ starting at $r_a$. These
projected paths cover $W_a$ and its edges, so $M_a$ is connected.
For an isolated target vertex it is the one-vertex graph.

\begin{lemma}\label{lem:reduction}
Suppose $X\subseteq E$ admits pairwise disjoint forests
$F_a\subseteq E_a\cap X$, each spanning every component of $M_a(X)$.
Then $G$ has a root-preserving minor $Q$ containing an $H$-scheme.
If $X\ne\varnothing$, then $|V(Q)|<|V(G)|$.
\end{lemma}

\begin{proof}
For every component $K$ of $M_a(X)$, set
\[
 D_{a,K}=V(K)\cup
 \{x\in F_a:\text{the edge labelled }x\text{ belongs to }K\}.
\]
Each allocated tree lifts to two-edge paths through its actual label
vertices, making $D_{a,K}$ connected. Isolated components give singleton
sets. These sets are pairwise disjoint: the original $A$ vertices belong
to distinct colour components, and every allocated $B$ label has just
one owner. They contain no $B$ root and at most one $A$ root.

Contract each $D_{a,K}$ to $d_{a,K}$, and delete every unallocated
vertex of $X$. Keep the roots $r_b$ and the vertices of $E\setminus X$
as separate vertices. Retain precisely the edges from $d_{a,K}$ to
surviving $B$ vertices $y$ witnessed by an original edge from $V(K)$
to $y$, deleting other edges and parallel copies. This defines a minor
$Q$ with a proper colouring: $d_{a,K}$ receives colour $a$, and the
surviving $B$ vertices retain their colours. The root of $a$ is its
component vertex, and every $B$ root is unchanged.

Fix $ab\in E(H)$. In the vertex sequence of $P_{ab}$, replace each
$A$ vertex by its component vertex, omit every $x\in X$, and suppress
consecutive repetitions. For each omitted segment $u x v$, the edge
labelled $x$ places $u,v$ in the same component of $M_a(X)$. Their
images therefore coincide, regardless of which forest owns $x$.
Every remaining step is witnessed by an original edge from an $A$
vertex to a surviving $B$ vertex. The sequence is consequently a walk
in $Q$ between the new roots, using only colours $a,b$.

Delete closed excursions to obtain a simple path. No foreign $A$ root
has colour $a$, and the only $B$ root retained from $P_{ab}$ is $r_b$.
Every intersection of the resulting paths has its colour as a common
target endpoint. Thus they form an $H$-scheme.

The quotient has $\sum_a c_a(X)+|B|+|E\setminus X|$ vertices, so
\[
 |V(G)|-|V(Q)|=|X|+\sum_a r_a(X).
\]
This is positive when $X\ne\varnothing$.
\end{proof}

\section{Proof of Theorem~\ref{thm:main}}

Remove isolated target vertices and their roots, restoring them as their
original singleton branch sets afterwards. An edgeless target is
immediate. For the remaining fixed target use strong induction on the
host order. Lemma~\ref{lem:normalization} gives a root-preserving minor
$K$ of the original host $G$, with a proper endpoint-coloured scheme and
$|V(K)|\leq|V(G)|$.

Orient the bipartition $(A,B)$ so that the number $N_A$ of nonroots
with colours in $A$ is at most the number $N_B$ with colours in $B$.
Form the projections in Section~3 for $K$. They satisfy
\[
 \sum_{a\in A}r_a(E)=\sum_{a\in A}(|W_a|-1)=N_A,
 \qquad |E|=N_B.
\]
Let $(I_a)$ be a maximising disjoint family of projection forests,
and let $R=\sum_a|I_a|$.

If $R=N_A$, every $I_a$ is a spanning tree. Set
\[
 C_a=W_a\cup I_a\quad(a\in A),\qquad C_b=\{r_b\}\quad(b\in B).
\]
The lifted trees make these sets connected, and the disjoint label sets
make them disjoint. They retain every root. The final edge of $P_{ab}$
joins $r_b$ to $W_a$, giving every required contact.

Otherwise $R<N_A\leq N_B=|E|$. A minimising set $X$ in
Lemma~\ref{lem:union} is nonempty, since the expression at the empty
set equals $|E|>R$. That lemma supplies the forests $F_a=I_a\cap X$
required by Lemma~\ref{lem:reduction}. We obtain a root-preserving minor
$Q$ containing an $H$-scheme, with
\[
 |V(Q)|<|V(K)|\leq |V(G)|.
\]
The induction hypothesis supplies a rooted model in $Q$. Replace its
vertices by their disjoint connected preimages, first in $K$ and then
in $G$. This preserves connectivity, disjointness, all roots and every
required adjacency. The same lift applies to the direct spanning-tree
case, completing the induction. \qed

The bipartition may reverse during induction, so lifted branch sets
may expand on both original shores.

\section{The bipartite-flow assertion}\label{sec:flow}

\begin{corollary}\label{cor:flow}
Let $H$ be bipartite of minimum degree at least two. Inject its vertices
into a finite graph $G$ as terminals and choose a simple path between
the corresponding terminal pair for every target edge. If paths for edges with four distinct endpoints are
vertex-disjoint, then $G$ has an $H$-minor rooted at every terminal.
\end{corollary}

\begin{proof}
A pairwise intersecting family of edges in a bipartite graph has a
common endpoint: after two edges $ab,ac$, an edge meeting both but
avoiding $a$ would be $bc$, forming a triangle. Hence paths meeting
at a vertex satisfy the common-endpoint condition. No terminal $r_v$
can be internal to a path for a nonincident edge $ab$: at most one of
$a,b$ is adjacent to $v$, so $v$ has a neighbour $w$ outside $\{a,b\}$.
Then $P_{vw}$ and $P_{ab}$ meet at $r_v$ although their edges have four
distinct endpoints. Thus the paths form a scheme, and
Theorem~\ref{thm:main} applies.
\end{proof}

Biswal, Lee and Rao~\cite[Lemma 3.2]{BLR} asserted this minor-existence
conclusion; their Lemma 3.4 explicitly seeks to retain every terminal.
The minimum-degree restriction does not distinguish their intended
rooted assertion from Theorem~\ref{thm:main}. To remove it, attach a
private four-cycle at every target vertex and a corresponding literal
cycle at its host root. The enlarged scheme has a bipartite target of
minimum degree two. All added host vertices are prescribed roots, so
none can belong to an original root's bag. Restricting a rooted model
to the original bags therefore proves contractibility of the original
target.

The published intersection definition~\cite[p.~13:10]{BLR} appears to
reverse the intended preprint convention: independent demand paths
must be vertex-disjoint. Correcting that wording does not repair the
prefix construction on p.~13:11. For example, in the coloured
$K_{2,2}$-scheme
\[
 P_{ij}=a_i\,b'_j\,a'_i\,b_j\qquad(i,j\in\{0,1\}),
\]
all eight displayed vertex names are distinct. Their left prefixes are
$\{a_i\}$, and their proposed right sets are
$\{b_j,b'_j,a'_0,a'_1\}$, which overlap. This refutes the construction's
Lemma~3.6, not the minor-existence conclusion: a rooted model is
$\{a_0,b'_0\}$, $\{a_1,b'_1\}$, $\{b_0,a'_1\}$,
$\{b_1,a'_0\}$, in the order $a_0,a_1,b_0,b_1$.

Lee~\cite[Lemma 3.3]{Lee} later invokes the same bipartite-flow
assertion. Some clique-flow applications now have another extraction
argument: Kolbe and Spalding-Jamieson~\cite[Lemma 3.11 and Proposition 3.4]{KSJ}, using
Korhonen and Lokshtanov~\cite[Lemmas 4.3--4.4]{KL}, choose a specific
bounded-degree demand graph for each clique order and extract an
unrooted clique minor. This does not establish contractibility for
arbitrary bipartite targets. Our proof supplies no new spectral,
separator or bounded-depth estimate.

\begin{thebibliography}{9}
\bibitem{BLR}
P.~Biswal, J.~R.~Lee and S.~Rao,
\href{https://doi.org/10.1145/1706591.1706593}{\emph{Eigenvalue bounds,
spectral partitioning, and metrical deformations via flows}},
J.~ACM \textbf{57}(3) (2010), Article 13.
\href{https://arxiv.org/abs/0808.0148v2}{arXiv:0808.0148v2}.

\bibitem{Edmonds}
J.~Edmonds,
\href{https://doi.org/10.1007/978-3-540-68279-0_7}{\emph{Matroid partition}},
reprinted in \emph{50 Years of Integer Programming 1958--2008},
Springer, 2010, 199--218.

\bibitem{KSJ}
B.~Kolbe and J.~Spalding-Jamieson,
\emph{On Eigenvalue Bounds for Bounded Genus Graphs and Minor-Free Graphs},
\href{https://arxiv.org/abs/2608.27179v1}{arXiv:2608.27179v1}, 2026.

\bibitem{KL}
T.~Korhonen and D.~Lokshtanov,
\href{https://doi.org/10.1137/1.9781611977912.188}{\emph{Induced-Minor-Free
Graphs: Separator Theorem, Subexponential Algorithms, and Improved
Hardness of Recognition}},
Proc.~SODA 2024, 5249--5275.
\href{https://arxiv.org/abs/2308.04795v1}{arXiv:2308.04795v1}.

\bibitem{KM}
M.~Kriesell and S.~Mohr, \emph{Kempe Chains and Rooted Minors},
\href{https://arxiv.org/abs/1911.09998v2}{arXiv:1911.09998v2}, 2022.

\bibitem{KPR}
A.~K\"undgen, M.~J.~Pelsmajer and R.~Ramamurthi,
\href{https://doi.org/10.1002/jgt.21812}{\emph{Finding minors in graphs
with a given path structure}},
J.~Graph Theory \textbf{79} (2015), 30--47.
\href{https://arxiv.org/abs/1207.6141}{arXiv:1207.6141}.

\bibitem{Lee}
J.~R.~Lee, \emph{Separators in region intersection graphs},
\href{https://arxiv.org/abs/1608.01612v3}{arXiv:1608.01612v3}, 2017.
\end{thebibliography}
\end{document}
